\documentclass[11pt, a4paper]{article}

\usepackage[a4paper, margin=0.9in]{geometry}
\usepackage[a4paper, margin=0.9in]{geometry}
\usepackage{cite}
\usepackage{amsmath,amssymb,amsfonts}
\usepackage{algorithmic}
\usepackage{graphicx}
\usepackage{textcomp}
\usepackage{xcolor}
\usepackage{needspace}
\usepackage{needspace}

\usepackage{fancyhdr}
\pagestyle{fancy}
\fancyhf{} 
\renewcommand{\headrulewidth}{0pt}
\cfoot{Page \thepage} 

\usepackage{hyperref}
\hypersetup{
    colorlinks=true,
    linkcolor=blue,
    urlcolor=blue,
    citecolor=blue
}

\begin{document}

\thispagestyle{fancy} 

\vspace*{-1cm}
\noindent\hspace*{-1.5cm}%
\begin{minipage}[c]{6.5cm}
    \raggedright
    \includegraphics[width=6.5cm]{figures/wits_logo.png}
\end{minipage}%
\hspace{0.5cm}
\begin{minipage}[c]{10.5cm}
    \centering
    {\LARGE \textbf{Assessment 1: Sourcing and Systematising Literature}\\[0.6cm]}
        {\Large Subverting Coding Agents with Poisoned Context\\[0.8cm]}
        {\Large Subverting Coding Agents with Poisoned Context\\[0.8cm]}
    {\large \textbf{Arno Strauss}\\[0.2cm]}
    {\normalsize \texttt{2613224@students.wits.ac.za}\\[0.1cm]}
    {\normalsize \selectfont April 10, 2025}
    {\normalsize \selectfont April 10, 2025}
\end{minipage}

\vspace{0cm}

\section*{References}

% PAPER 1 (D)
\noindent \hypertarget{p1}{[Debenedetti et al. 2024]} E. Debenedetti, J. Zhang, M. Balunovic, L. Beurer-Kellner, M. Fischer, and F. Tram\`er. \href{https://arxiv.org/abs/2406.13352}{AgentDojo: A Dynamic Environment to Evaluate Prompt Injection Attacks and Defenses for LLM Agents.} \textit{Proceedings of the 38th Conference on Neural Information Processing Systems (NeurIPS), 2024.}
\begin{quote}
\textbf{Aim:} The authors set out to construct a multi-turn, stateful evaluation framework that specifically measures how tool-calling agents respond to adversarial prompt injections.
\textbf{Aim:} The authors set out to construct a multi-turn, stateful evaluation framework that specifically measures how tool-calling agents respond to adversarial prompt injections.

\vspace{0.5em}
\textbf{Style/Type:} Conference paper, empirical benchmark.
\textbf{Style/Type:} Conference paper, empirical benchmark.

\vspace{0.5em}
\textbf{Summary:} This research addresses the instability in current security conclusions caused by the overlap between model capabilities and evaluator limitations. By introducing Agent Dojo, the researchers create a testing environment that verifies outcomes through actual state changes rather than simple text matching. Their findings highlight a concerning trend where increased model intelligence actually correlates with higher vulnerability to instruction-precedence attacks. \\ \textit{Research question} Under what conditions do deterministic state metrics capture security-utility trade-offs, and do defenses reduce ASR without unacceptable utility loss? \\ \textit{Key innovation} The method decouples attack verification from model self-evaluation by checking environment state. \\ \textit{Boundary Conditions} The interaction scope is limited strictly to text, which excludes multimodal attack surfaces. Furthermore, task definitions, tool schemas, and utility checks rely on manual configuration, limiting overall scalability.
\textbf{Summary:} This research addresses the instability in current security conclusions caused by the overlap between model capabilities and evaluator limitations. By introducing Agent Dojo, the researchers create a testing environment that verifies outcomes through actual state changes rather than simple text matching. Their findings highlight a concerning trend where increased model intelligence actually correlates with higher vulnerability to instruction-precedence attacks. \\ \textit{Research question} Under what conditions do deterministic state metrics capture security-utility trade-offs, and do defenses reduce ASR without unacceptable utility loss? \\ \textit{Key innovation} The method decouples attack verification from model self-evaluation by checking environment state. \\ \textit{Boundary Conditions} The interaction scope is limited strictly to text, which excludes multimodal attack surfaces. Furthermore, task definitions, tool schemas, and utility checks rely on manual configuration, limiting overall scalability.

\vspace{0.5em}
\Needspace{8\baselineskip}
\textbf{Cross-references:} This paper provides the adversarial measurement layer for this project. \texttt{SWE-bench} in \hyperlink{p2}{[Jimenez et al. 2024]} provides the baseline for functional repair utility, while AgentDojo provides the adversarial metrics needed to evaluate the security attacks shown in \hyperlink{p4}{[Przymus et al. 2026]}. It evaluates the agent design in \hyperlink{p7}{[Yang et al. 2024]}, tests attacks from \hyperlink{p3}{[Liu et al. 2025]} and \hyperlink{p4}{[Przymus et al. 2026]}, and measures the effect of \texttt{PromptArmor} in \hyperlink{p5}{[Shi et al. 2025]}.

\vspace{0.5em}
\Needspace{4\baselineskip}
\enlargethispage{0.8\baselineskip}
\textbf{Relevance:} This entry defines how the project measures security and utility together. A critical tension emerges here as better tool competence simultaneously raises benign utility and attack success. This means capability growth inadvertently increases risk within the same operating regime.
\Needspace{4\baselineskip}
\enlargethispage{0.8\baselineskip}
\textbf{Relevance:} This entry defines how the project measures security and utility together. A critical tension emerges here as better tool competence simultaneously raises benign utility and attack success. This means capability growth inadvertently increases risk within the same operating regime.
\end{quote}

\newpage
% PAPER 2 (J)
\noindent \hypertarget{p2}{[Jimenez et al. 2024]} C. E. Jimenez et al. \href{https://arxiv.org/abs/2310.06770}{\texttt{SWE-bench}: Can language models resolve real-world GitHub issues?} \textit{Proceedings of the International Conference on Learning Representations (ICLR)}, 2024.
\noindent \hypertarget{p2}{[Jimenez et al. 2024]} C. E. Jimenez et al. \href{https://arxiv.org/abs/2310.06770}{\texttt{SWE-bench}: Can language models resolve real-world GitHub issues?} \textit{Proceedings of the International Conference on Learning Representations (ICLR)}, 2024.
\begin{quote}
\textbf{Aim:} This study introduces a repository-scale benchmark that evaluates how well language models can resolve complex, real-world software engineering tasks across large codebases.
\textbf{Aim:} This study introduces a repository-scale benchmark that evaluates how well language models can resolve complex, real-world software engineering tasks across large codebases.

\vspace{0.5em}
\textbf{Style/Type:} Conference paper, empirical benchmark.

\vspace{0.5em}
\textbf{Summary:} SWE-bench focuses on the difficulties of distinguishing between an agent’s ability to locate a bug and its ability to actually fix it. Using a rigorous three-stage curation pipeline, the authors demonstrate that even advanced models struggle with context overload. They find that models often prioritize 'greedy' patches that solve immediate tests but degrade the long-term maintainability of the software. \\ \textit{Research question} Whether execution-based benchmarking separates repository reasoning across BM25 and Oracle retrieval, and whether pass outcomes track maintainability quality? \\ \textit{Key innovation} The framework decouples curation, execution validation, and quality assessment. This allows separate measurement of resolution, regression, and structural patch complexity. \\ \textit{Boundary Conditions} Generalizability is inherently limited by the benchmark's Python-only scope, which restricts transferability to polyglot repositories. Additionally, the evaluation focuses on non-agentic large language model baselines rather than fully autonomous agents.
\textbf{Summary:} SWE-bench focuses on the difficulties of distinguishing between an agent’s ability to locate a bug and its ability to actually fix it. Using a rigorous three-stage curation pipeline, the authors demonstrate that even advanced models struggle with context overload. They find that models often prioritize 'greedy' patches that solve immediate tests but degrade the long-term maintainability of the software. \\ \textit{Research question} Whether execution-based benchmarking separates repository reasoning across BM25 and Oracle retrieval, and whether pass outcomes track maintainability quality? \\ \textit{Key innovation} The framework decouples curation, execution validation, and quality assessment. This allows separate measurement of resolution, regression, and structural patch complexity. \\ \textit{Boundary Conditions} Generalizability is inherently limited by the benchmark's Python-only scope, which restricts transferability to polyglot repositories. Additionally, the evaluation focuses on non-agentic large language model baselines rather than fully autonomous agents.

\vspace{0.5em}
\Needspace{8\baselineskip}
\textbf{Cross-references:} \texttt{SWE-bench} is the execution baseline for this project; on top of it, \hyperlink{p7}{[Yang et al. 2024]} supplies the foundational agent design, and \hyperlink{p4}{[Przymus et al. 2026]} shows how issue text weaponizes that workflow. \texttt{PromptArmor} in \hyperlink{p5}{[Shi et al. 2025]} is the corresponding defense, while \hyperlink{p3}{[Liu et al. 2025]} adds host-level injection paths beyond repository scope.

\vspace{0.5em}
\Needspace{8\baselineskip}
\textbf{Relevance:} This entry provides the baseline needed to measure attack and defense effects. This reveals a structural blind spot where high pass-to-pass utility frequently coexists with compromised security and maintainability. Consequently, tests can stay green while exploit-preserving repairs easily pass through.
\Needspace{8\baselineskip}
\textbf{Relevance:} This entry provides the baseline needed to measure attack and defense effects. This reveals a structural blind spot where high pass-to-pass utility frequently coexists with compromised security and maintainability. Consequently, tests can stay green while exploit-preserving repairs easily pass through.
\end{quote}

\newpage
% PAPER 3 (L)
\noindent \hypertarget{p3}{[Liu et al. 2025]} Y. Liu, Y. Zhao, Y. Lyu, T. Zhang, H. Wang, and D. Lo. \href{https://arxiv.org/abs/2509.22040}{``Your AI, My Shell'': Demystifying Prompt Injection Attacks on Agentic AI Coding Editors.} \textit{Proceedings of the ACM on Software Engineering (PACMSE)}, 2025.
\begin{quote}
\textbf{Aim:} The researchers investigate the specific security risks associated with high-privilege agentic coding editors when they are exposed to untrusted instruction channels.
\textbf{Aim:} The researchers investigate the specific security risks associated with high-privilege agentic coding editors when they are exposed to untrusted instruction channels.

\vspace{0.5em}
\textbf{Style/Type:} Journal article, empirical study.
\textbf{Style/Type:} Journal article, empirical study.

\vspace{0.5em}
\textbf{Summary:} This study demonstrates how policy controls in modern editors can be bypassed by injecting malicious logic into seemingly benign context. Through the AIShell Jack framework, the authors show that autonomous modes are particularly susceptible to unauthorized actions. They use semantic intent matching to prove that advanced reasoning actually makes it easier for an agent to execute a harmful exploit efficiently. \\ \textit{Research question} Which configuration conditions increase prompt-injection risk in high-privilege editors, and do unauthorized-action classes differ in attack rates? \\ \textit{Key innovation} The framework decouples success scoring from literal command matching by using a semantic-intent criterion. \\ \textit{Boundary Conditions} Relying on payloads without advanced obfuscation may understate true adversary capabilities. Furthermore, automated analysis primarily targets revised code and executed commands since conversational logs remain hidden behind commercial editor isolation policies.
\textbf{Summary:} This study demonstrates how policy controls in modern editors can be bypassed by injecting malicious logic into seemingly benign context. Through the AIShell Jack framework, the authors show that autonomous modes are particularly susceptible to unauthorized actions. They use semantic intent matching to prove that advanced reasoning actually makes it easier for an agent to execute a harmful exploit efficiently. \\ \textit{Research question} Which configuration conditions increase prompt-injection risk in high-privilege editors, and do unauthorized-action classes differ in attack rates? \\ \textit{Key innovation} The framework decouples success scoring from literal command matching by using a semantic-intent criterion. \\ \textit{Boundary Conditions} Relying on payloads without advanced obfuscation may understate true adversary capabilities. Furthermore, automated analysis primarily targets revised code and executed commands since conversational logs remain hidden behind commercial editor isolation policies.

\vspace{0.5em}
\Needspace{8\baselineskip}
\textbf{Cross-references:} This study explains host-level attack paths that complement repository-level attacks in \hyperlink{p4}{[Przymus et al. 2026]}, while the victim interface comes from \hyperlink{p7}{[Yang et al. 2024]}. The adoption baseline in \hyperlink{p6}{[Watanabe et al. 2025]} explains exposure frequency, and \texttt{PromptArmor} in \hyperlink{p5}{[Shi et al. 2025]} provides the matching pre-execution defense.

\vspace{0.5em}
\Needspace{8\baselineskip}
\textbf{Relevance:} This entry strengthens the external-context threat model with host-level evidence. An unintended consequence of advanced models is that stronger reasoning directly increases exploit execution efficiency. As a result, capability gains actively amplify security loss when untrusted instruction channels remain open.
\Needspace{8\baselineskip}
\textbf{Relevance:} This entry strengthens the external-context threat model with host-level evidence. An unintended consequence of advanced models is that stronger reasoning directly increases exploit execution efficiency. As a result, capability gains actively amplify security loss when untrusted instruction channels remain open.
\end{quote}

\newpage
% PAPER 4 (P)
\noindent \hypertarget{p4}{[Przymus et al. 2026]} P. Przymus, A. Happe, and J. Cito. \href{https://arxiv.org/abs/2509.05372}{Adversarial Bug Reports as a Security Risk in Language Model-Based Automated Program Repair.} In \textit{Proceedings of the 23rd International Conference on Mining Software Repositories (MSR)}, 2026.
\begin{quote}
\textbf{Aim:} This paper evaluates whether adversarial bug reports can be weaponized to force automated program repair systems into generating insecure or harmful code.
\textbf{Aim:} This paper evaluates whether adversarial bug reports can be weaponized to force automated program repair systems into generating insecure or harmful code.

\vspace{0.5em}
\textbf{Style/Type:} Conference paper, empirical study.
\textbf{Style/Type:} Conference paper, empirical study.

\vspace{0.5em}
\textbf{Summary:} The authors highlight a massive economic disparity between the negligible cost of generating an attack and the high cost of executing the repair pipeline. By mapping red-teaming efforts to STRIDE tactics, they show that most malicious patches successfully pass standard continuous integration checks. This suggests that current security filters are largely ineffective against narratives that hide exploits within plausible-sounding reports. \\ \textit{Research question} If adversarial bug reports are introduced, do APR workflows produce insecure yet test-passing outputs, and does mitigation effectiveness vary by workflow stage? \\ \textit{Key innovation} The framework decouples adversarial-report synthesis from repair execution through context-aware red-teaming mapped to STRIDE tactics. \\ \textit{Boundary Conditions} The scope of these findings is restricted to a single automated program repair tool, leaving architectural variations untested. Validating attack success also relies heavily on manual expert review, which introduces inherent subjectivity risks.
\textbf{Summary:} The authors highlight a massive economic disparity between the negligible cost of generating an attack and the high cost of executing the repair pipeline. By mapping red-teaming efforts to STRIDE tactics, they show that most malicious patches successfully pass standard continuous integration checks. This suggests that current security filters are largely ineffective against narratives that hide exploits within plausible-sounding reports. \\ \textit{Research question} If adversarial bug reports are introduced, do APR workflows produce insecure yet test-passing outputs, and does mitigation effectiveness vary by workflow stage? \\ \textit{Key innovation} The framework decouples adversarial-report synthesis from repair execution through context-aware red-teaming mapped to STRIDE tactics. \\ \textit{Boundary Conditions} The scope of these findings is restricted to a single automated program repair tool, leaving architectural variations untested. Validating attack success also relies heavily on manual expert review, which introduces inherent subjectivity risks.

\vspace{0.5em}
\Needspace{8\baselineskip}
\textbf{Cross-references:} This entry defines the attack path and targets the victim architecture formalized within the \texttt{SWE-agent} framework \hyperlink{p7}{[Yang et al. 2024]}. It extends the benign execution baseline provided by the \texttt{SWE-bench} dataset \hyperlink{p2}{[Jimenez et al. 2024]}, aligns with host-level paths in \hyperlink{p3}{[Liu et al. 2025]}, and connects directly to the defense in \texttt{PromptArmor} \hyperlink{p5}{[Shi et al. 2025]}.

\vspace{0.5em}
\Needspace{8\baselineskip}
\textbf{Relevance:} This entry is the direct empirical core of the project threat model. This exposes a severe economic asymmetry because low-cost attack synthesis at approximately \$0.0001 drives high-cost processing near \$100. Compounding this issue, 68.6\% of malicious outputs still pass CI, which effectively masks security failures behind functional utility.
\Needspace{8\baselineskip}
\textbf{Relevance:} This entry is the direct empirical core of the project threat model. This exposes a severe economic asymmetry because low-cost attack synthesis at approximately \$0.0001 drives high-cost processing near \$100. Compounding this issue, 68.6\% of malicious outputs still pass CI, which effectively masks security failures behind functional utility.
\end{quote}

\newpage
% PAPER 5 (S)
\noindent \hypertarget{p5}{[Shi et al. 2025]} T. Shi et al. \href{https://arxiv.org/abs/2507.15219}{\texttt{PromptArmor}: Simple yet Effective Prompt Injection Defenses.} \textit{arXiv preprint arXiv:2507.15219}, 2025.
% PAPER 5 (S)
\noindent \hypertarget{p5}{[Shi et al. 2025]} T. Shi et al. \href{https://arxiv.org/abs/2507.15219}{\texttt{PromptArmor}: Simple yet Effective Prompt Injection Defenses.} \textit{arXiv preprint arXiv:2507.15219}, 2025.
\begin{quote}
\textbf{Aim:} The study presents a preprocessing guardrail intended to sanitize external data before it can be consumed by tool-using agents.
\textbf{Aim:} The study presents a preprocessing guardrail intended to sanitize external data before it can be consumed by tool-using agents.

\vspace{0.5em}
\textbf{Style/Type:} Preprint, empirical study.

\vspace{0.5em}
\textbf{Summary:} PromptArmor functions as a protective wrapper that processes untrusted inputs using a secondary guardrail model. The researchers found that while this method effectively neutralizes most injection attacks on high-capacity models, its effectiveness drops significantly on smaller architectures. This suggests that the security of such wrappers is currently tied directly to the available compute power. \\ \textit{Research question} Can off-the-shelf LLM sanitization keep ASR below $1\%$ while preserving agent task-utility thresholds under adaptive attack pressure? \\ \textit{Key innovation} The approach decouples untrusted-input sanitization from downstream task reasoning through a training-free wrapper path. \\ \textit{Boundary Conditions} Defensive efficacy is heavily dependent on model capacity. Smaller models exhibit severe security-utility trade-offs, and high prompt sensitivity means robust behavior requires explicit prompt-injection definitions rather than generic safety guidelines.
\textbf{Summary:} PromptArmor functions as a protective wrapper that processes untrusted inputs using a secondary guardrail model. The researchers found that while this method effectively neutralizes most injection attacks on high-capacity models, its effectiveness drops significantly on smaller architectures. This suggests that the security of such wrappers is currently tied directly to the available compute power. \\ \textit{Research question} Can off-the-shelf LLM sanitization keep ASR below $1\%$ while preserving agent task-utility thresholds under adaptive attack pressure? \\ \textit{Key innovation} The approach decouples untrusted-input sanitization from downstream task reasoning through a training-free wrapper path. \\ \textit{Boundary Conditions} Defensive efficacy is heavily dependent on model capacity. Smaller models exhibit severe security-utility trade-offs, and high prompt sensitivity means robust behavior requires explicit prompt-injection definitions rather than generic safety guidelines.

\vspace{0.5em}
\Needspace{8\baselineskip}
\textbf{Cross-references:} This paper is the defense against attack pathways in \hyperlink{p4}{[Przymus et al. 2026]} and \hyperlink{p3}{[Liu et al. 2025]}. It applies to the victim surface defined by \hyperlink{p7}{[Yang et al. 2024]}, and it is tested on the execution baseline in \hyperlink{p2}{[Jimenez et al. 2024]}.

\vspace{0.5em}
\Needspace{8\baselineskip}
\textbf{Relevance:} This entry provides the main mitigation candidate for the project threat model. A notable trade-off in this defense is its scale dependency. While attack neutralization and utility preservation coexist seamlessly at high-capacity settings, the design degrades rapidly on smaller models to bind security gains strictly to compute availability.
\Needspace{8\baselineskip}
\textbf{Relevance:} This entry provides the main mitigation candidate for the project threat model. A notable trade-off in this defense is its scale dependency. While attack neutralization and utility preservation coexist seamlessly at high-capacity settings, the design degrades rapidly on smaller models to bind security gains strictly to compute availability.
\end{quote}

\newpage
% PAPER 6 (A)
\noindent \hypertarget{p6}{[Watanabe et al. 2025]} M. Watanabe, H. Li, Y. Kashiwa, B. Reid, H. Iida, and A. E. Hassan. \href{https://arxiv.org/abs/2509.14745}{On the Use of Agentic Coding: An Empirical Study of Pull Requests on GitHub.} \textit{ACM Transactions on Software Engineering and Methodology (TOSEM)}, 2025.
\begin{quote}
\textbf{Aim:} This empirical study examines the real-world adoption and acceptance rates of agent-generated pull requests compared to those created by human developers.

\vspace{0.5em}
\textbf{Style/Type:} Journal article, empirical study.

\vspace{0.5em}
\textbf{Summary:} By analyzing hundreds of matched pull requests on GitHub, the authors find that agent-generated code is accepted at a high rate but often receives less rigorous scrutiny. The study introduces 'Confidence Cards' as a way to improve the transparency of agentic plans. Their data suggests that the speed of automated merging may be masking significant risks regarding the depth of code reviews. \\ \textit{Research question} To what extent do Agentic-PRs meet acceptance and zero-revision thresholds versus Human-PRs, and do scope differences predict review friction? \\ \textit{Key innovation} The study decouples acceptance outcomes from informal review impressions through matched Agentic-PR and Human-PR analysis. It also adds a ``Confidence Card'' to separate plan traceability from diff inspection. \\ \textit{Boundary Conditions} Relying on explicit text markers fails to capture hybrid workflows where developers manually modify agentic code before submission. Moreover, early adoption bias during the studied period might artificially inflate acceptance signals.

\vspace{0.5em}
\Needspace{8\baselineskip}
\textbf{Cross-references:} This entry shows how threats appear in production pull requests. \hyperlink{p7}{[Yang et al. 2024]} provides the foundational model of agent behavior, and \hyperlink{p4}{[Przymus et al. 2026]} shows how adversarial bug reports exploit that behavior. \hyperlink{p3}{[Liu et al. 2025]} aligns with MCP and shell-directed manipulation patterns, while \hyperlink{p1}{[Debenedetti et al. 2024]} provides security-utility metrics for measurement. \texttt{PromptArmor} in \hyperlink{p5}{[Shi et al. 2025]} is the matching defense under reviewer workload limits.

\vspace{0.5em}
\Needspace{8\baselineskip}
\textbf{Relevance:} This entry grounds the threat model in real repositories rather than only benchmarks. The central vulnerability in production is speed. While 54.9\% zero-revision merging accelerates delivery, it severely reduces inspection depth and actively elevates the risk of adversarial insertions via manipulated external context.
\end{quote}


\newpage
% PAPER 7 (Y)
\noindent \hypertarget{p7}{[Yang et al. 2024]} J. Yang, C. E. Jimenez, A. Wettig, K. Lieret, S. Yao, K. Narasimhan, and O. Press. \href{https://arxiv.org/abs/2405.15793}{\texttt{SWE-agent}: Agent-Computer Interfaces Enable Automated Software Engineering.} \textit{Proceedings of the 38th Conference on Neural Information Processing Systems (NeurIPS)}, 2024.
\begin{quote}
\textbf{Aim:} The authors propose a specialized interface designed specifically for agent constraints to improve performance on automated software engineering tasks.
\textbf{Aim:} The authors propose a specialized interface designed specifically for agent constraints to improve performance on automated software engineering tasks.

\vspace{0.5em}
\textbf{Style/Type:} Conference paper, empirical study and system architecture.
\textbf{Style/Type:} Conference paper, empirical study and system architecture.

\vspace{0.5em}
\textbf{Summary:} Instead of forcing agents to use human-centric tools, SWE-agent utilizes a tailored Agent-Computer Interface that limits search results and enforces syntax checks. While this architecture significantly improves the success rate of bug resolutions, it also creates a highly predictable pathway for attackers. The interface structure itself becomes the primary variable that determines both the reliability and the vulnerability of the system. \\ \textit{Research question} To what extent does replacing human-oriented interfaces with an Agent-Computer Interface improve \texttt{SWE-bench} resolution and cost efficiency versus non-interactive RAG? \\ \textit{Key innovation} The design decouples agent interaction from human shell and IDE affordances. It separates retrieval control, local file viewing, and syntax-constrained editing recovery. \\ \textit{Boundary Conditions} The current toolkit remains basic and excludes advanced software engineering practices like fuzz testing or web browsing. Because interface components and recovery guardrails are manually engineered from observed failure modes, automatic optimization and dynamic interface evolution are not yet supported. This implies interface structure acts as a primary control variable for agent reliability.

\vspace{0.5em}
\Needspace{8\baselineskip}
\textbf{Cross-references:} This entry defines the victim architecture on the \texttt{SWE-bench} baseline \hyperlink{p2}{[Jimenez et al. 2024]}. The adversarial bug-report path in \hyperlink{p4}{[Przymus et al. 2026]} targets this interface and control loop, while \hyperlink{p3}{[Liu et al. 2025]} extends the attack path at host level and \hyperlink{p5}{[Shi et al. 2025]} provides the defense option.

\vspace{0.5em}
\Needspace{8\baselineskip}
\textbf{Relevance:} This entry specifies the victim mechanics for the project threat model. This architecture presents a double-edged sword because the specific interface constraints that enforce execution discipline and improve repair utility simultaneously map out a highly predictable ingestion pathway for adversarial context.
\end{quote}


\end{document}